\documentclass[11pt]{article}
\usepackage[utf8]{inputenc}
\usepackage[T1]{fontenc}
\usepackage[margin=1in]{geometry}
\usepackage{amsmath,amssymb}
\usepackage{natbib}
\usepackage{graphicx}
\graphicspath{{figures/}}
\usepackage{booktabs}
\usepackage{caption}
\usepackage[hidelinks]{hyperref}
\captionsetup{font=small,labelfont=bf}
\setlength{\parskip}{2pt}
\setlength{\emergencystretch}{3em}

\title{\bfseries Regular spacing of salt structures as a Rayleigh--Taylor signature:
edge-corrected point-pattern and structure-factor analysis of the U.S.\ Gulf Coast and North German basins}
\author{R.~Chen\\{\small GUT Geoservice Inc., Montr\'eal, Qu\'ebec, Canada}}
\date{\today}

\begin{document}
\maketitle

\begin{abstract}
\noindent
Salt structures rise by a buoyancy-driven (Rayleigh--Taylor) instability that, in theory, selects a
characteristic wavelength and should therefore imprint a \emph{regular} map-view spacing of diapirs and
pillows. Empirically the record is mixed: several provinces show little or no characteristic wavelength,
and Rayleigh--Taylor models often fail to reproduce observed spacing. Whether complete salt-structure
\emph{populations} are statistically regular in map view has rarely been tested with modern edge-corrected
spatial statistics. We analyse two large open datasets --- 1253 Gulf Coast salt diapirs (USGS) and 693
North German salt structures (BGR InSpEE) --- using the Clark--Evans nearest-neighbour index, the pair
correlation function $g(r)$, Ripley's $L$, and the structure factor $S(k)$ with number variance, against
both complete-spatial-randomness (CSR) and inhomogeneous-Poisson null models. Pooled across a whole basin
the structures are \emph{clustered}, but this reflects the superposition of distinct sub-basins; within a
single coherent loading-driven province the structures are significantly \emph{regularly spaced}
(Clark--Evans $R=1.10$--$1.47$, $p\approx0.002$; robust to a density-gradient inhomogeneous-Poisson null),
with hard-core exclusion and a 12--15\,km characteristic spacing. Regularity strengthens with salt
thickness (Gl\"uckstadt Graben $>$ central basin $>$ Gulf shelf), as expected for wavelength selection. The
mature-diapir sub-class \emph{alone} is clustered because mature diapirs concentrate in depocentres, so
sub-class selection must be avoided. The structure factor shows genuine intermediate-range suppression
($S\!\approx\!0.42$ at a $\sim$37\,km scale, below the gradient-matched null) but the long-wavelength limit
is masked by the intrinsic basin density gradient, so strong hyperuniformity cannot be confirmed. The
results provide a quantitative, statistically hardened map-view signature of Rayleigh--Taylor wavelength
selection, together with explicit cautions on superposition, sub-class selection, and density gradients.
\end{abstract}

\section{Introduction}
Salt flows on geological timescales and, being less dense than its overburden, is gravitationally unstable:
a buried salt layer overlain by denser sediment is subject to a Rayleigh--Taylor (RT) instability whose
linear theory selects a dominant wavelength set primarily by the salt-layer thickness and the
viscosity/strength contrast \citep{trusheim1960,hudec2007,fernandez2015}. Because the fastest-growing mode
has a finite wavelength, the classical expectation is that salt pillows and the diapirs that evolve from
them should be \emph{quasi-regularly spaced} in map view, each structure surrounded by a salt-withdrawal
rim syncline that suppresses the growth of close neighbours.

Whether real provinces realise this prediction is, however, unsettled. Spacing/thickness relations are
reproduced in some settings but not others; variogram and spectral analyses of certain provinces report
essentially stochastic salt-surface topography with no uniform wavelength, and detailed studies find that
simple RT models often fail to reproduce observed diapir spacing once sedimentation, basement structure and
later shortening are accounted for \citep{jackson2017}. Real salt tectonics is further complicated by
structural inheritance: where salt ascent is localised by pre-existing faults, structures tend to align and
\emph{cluster} rather than disperse. The map-view spacing of a complete salt-structure population is
therefore an open empirical question, and one that is well posed for spatial point-pattern analysis.

Surprisingly, although nearest-neighbour and spacing statistics have long been used on individual salt
provinces, complete salt-structure populations have rarely been examined with modern edge-corrected
point-pattern statistics combined with the structure factor, explicit null models, and attention to the
basin-scale density gradient that is intrinsic to salt provinces. Here we do so for two large, openly
licensed datasets that together span the two end-member tectonic settings of passive-margin
sediment-loading (U.S.\ Gulf Coast) and an intracratonic basin with both loading- and inversion-driven salt
movement (North German Basin). We ask three questions: (i) are loading-driven salt-structure populations
significantly regularly spaced, and is any regularity robust to the basin density gradient; (ii) does the
strength of regularity scale with salt thickness, as RT wavelength selection predicts; and (iii) what is
the \emph{character} of the order --- merely a liquid-like spacing, or a longer-range suppression of density
fluctuations?

\section{Data}
\paragraph{U.S.\ Gulf Coast (USGS).}
We use the U.S.\ Geological Survey compilation \emph{Salt Diapirs in the Gulf Coast} (gcdiapirg;
\citealp{usgsgc}), itself derived from \citet{martin1980} and regional structural compilations. The dataset
contains 1253 salt-diapir outline polygons (NAD83 geographic coordinates) spanning the onshore and offshore
Gulf Coast from Texas to Florida and south to the Bay of Campeche. We take the area-weighted polygon
centroid as the point representing each structure and project geographic coordinates to a local Cartesian
frame (equirectangular projection at the mean latitude); all distances are reported in kilometres.

\paragraph{North German Basin (BGR InSpEE).}
We use the BGR \emph{InSpEE-Salzstrukturen} dataset (\citealp{bgr}; base mapping after
\citealp{reinhold2008}), an open shapefile (ETRS89 / UTM zone 32N, metres) of salt-structure outlines in
Northern Germany. The 855 polygons resolve into 693 named structures (area-weighted centroid per structure).
A companion internal-structure-type layer classifies each structure as a salt diapir
(\textit{Salzdiapir}, $n=159$), a compressively overprinted diapir ($n=107$), a salt pillow
(\textit{Salzkissen}, $n=205$), a double-salinar structure ($n=100$), a tectonic structure ($n=86$), or
unassigned. We define the \emph{halokinetic} population as the loading/ascent-driven structures (diapirs,
compressively overprinted diapirs, pillows and double-salinar structures; $n=571$), excluding the purely
tectonic and unassigned ones.

\section{Methods}
\paragraph{Province subsetting.}
A point pattern pooled over several distinct sub-basins is the superposition of independent processes and
will generally appear clustered even if each sub-basin is internally regular. We therefore report whole-basin
statistics for context but base inference on coherent single-province sub-regions of approximately uniform
tectonic setting.

\paragraph{Nearest-neighbour index.}
The Clark--Evans index \citep{clark1954} is $R=\bar d_{\mathrm{NN}}/\bar d_{\mathrm{CSR}}$, where
$\bar d_{\mathrm{NN}}$ is the mean nearest-neighbour distance. To remove edge bias we estimate
$\bar d_{\mathrm{CSR}}$ by Monte-Carlo simulation of complete spatial randomness (CSR) in the identical
convex-hull window (here $S=499$ realisations), so that $R>1$ indicates dispersion/regularity and $R<1$
clustering relative to the matched null; a one-sided $p$-value for dispersion is obtained by rank.

\paragraph{Second-order statistics.}
The pair correlation function $g(r)$ and Ripley's $K$ (reported as $L(r)-r=\sqrt{K(r)/\pi}-r$) are computed
with the border (reduced-sample) edge correction \citep{ripley1977,baddeley2015}. $g(r)=0$ at small $r$
measures hard-core exclusion; the first $g(r)$ maximum gives the characteristic spacing; $L(r)-r<0$
indicates regularity. Significance is assessed with the global studentised maximum-deviation envelope test
\citep{myllymaki2017} over the interaction range, against (i) a CSR null and (ii) an inhomogeneous-Poisson
null whose intensity is a Gaussian kernel estimate of the observed intensity (bandwidth $2\times$ median
nearest-neighbour distance). The inhomogeneous null builds the basin density gradient into the reference, so
that a significant result is not an artifact of the gradient; because a smooth gradient biases the
statistics toward \emph{clustering}, detection of regularity against this null is conservative.

\paragraph{Structure factor and number variance.}
The structure factor $S(\mathbf k)=N^{-1}\big|\sum_j e^{-i\mathbf k\cdot\mathbf r_j}\big|^2$ is evaluated at
the wavevectors allowed by the rectangular bounding window and radially averaged; this direct estimator
avoids density-field gridding artifacts. We also compute the number variance $\sigma^2(R)$ of point counts
in interior discs of radius $R$. Both are compared to CSR and inhomogeneous-Poisson nulls. A long-wavelength
limit $S(k\!\to\!0)\to0$ (equivalently $\sigma^2(R)$ growing more slowly than the window area) would indicate
hyperuniformity \citep{torquato2018}; a finite $S(k\!\to\!0)$ with a packing peak indicates a liquid-like
regular pattern.

\section{Results}

\subsection{Gulf Coast: clustered when pooled, regular within provinces}
Pooled over the whole Gulf the diapirs are strongly clustered (Clark--Evans $R=0.65$; Fig.~\ref{fig:gulfmap}),
as expected for the superposition of several salt basins separated by salt-free ground. Restricting to single
coherent loading-driven provinces reverses the verdict: the Texas--Louisiana shelf core ($N=271$) and the
broader northern-Gulf shelf/slope and coastal province ($N=571$) are significantly \emph{regular}
($R=1.12$ and $1.10$; one-sided $p=0.002$; Table~\ref{tab:R}). The Bay of Campeche is indistinguishable from
random, and the onshore interior salt basins, pooled across separate depocentres, are clustered
(Table~\ref{tab:R}).

The pair correlation of the Texas--Louisiana province (Fig.~\ref{fig:gulfgr}) shows hard-core exclusion
($g(r)=0$ below $\sim$4\,km), a clear packing peak ($g\!\approx\!1.5$) near 13\,km, and damped oscillations
toward unity --- the classic signature of a regular/liquid-like pattern. The observed $g(r)$ lies below both
the CSR and inhomogeneous-Poisson envelopes at small $r$; the global envelope test over $r\le16$\,km gives
$p=0.013$ against CSR and $p=0.010$ against the inhomogeneous null. Correspondingly $L(r)-r<0$ for
$r\lesssim13$\,km. (At larger $r$, $L(r)-r$ turns positive; for this elongate shelf band this is a
convex-hull window artifact rather than true clustering, and is not interpreted as long-range structure.)

\begin{figure}[t]
\centering
\includegraphics[width=\textwidth]{gulf_overview.pdf}
\caption{Gulf Coast salt diapirs (USGS gcdiapirg, $N=1253$). Left: centroids of diapir polygons. Right:
binned counts highlighting the separate salt basins (northern-Gulf shelf/slope and coastal trend; Bay of
Campeche) and the strong basin-scale density variation. Pooled over all basins the pattern is clustered
($R=0.65$); single coherent provinces are regular (Table~\ref{tab:R}).}
\label{fig:gulfmap}
\end{figure}

\begin{figure}[t]
\centering
\includegraphics[width=0.74\textwidth]{gulf_pair_correlation.pdf}
\caption{Pair correlation $g(r)$ of the Texas--Louisiana shelf province with $95\%$ Monte-Carlo envelopes
for CSR (grey) and an inhomogeneous-Poisson null carrying the density gradient (orange). Hard-core
exclusion below $\sim$4\,km and a packing peak near 13\,km; the observed curve falls below both envelopes at
small $r$ (global test $p=0.013$ vs.\ CSR, $p=0.010$ vs.\ inhomogeneous null).}
\label{fig:gulfgr}
\end{figure}

\subsection{North German Basin: stronger, cleaner regularity}
The North German salt structures are significantly regular already as a whole population ($R=1.22$, $N=693$;
$R=1.21$ for the halokinetic subset). Restricting to the central basin sharpens the signal markedly: the
central-basin halokinetic structures ($N=302$) give $R=1.41$ ($p=0.002$), and the Gl\"uckstadt Graben ---
the deepest, thickest-salt depocentre --- gives the strongest value, $R=1.47$ (Table~\ref{tab:R};
Fig.~\ref{fig:degrmap}).

The second-order statistics (Fig.~\ref{fig:degrgr}) confirm a strong, clean regularity: $g(r)=0$ below
$\sim$7\,km, a packing peak near 15\,km, and $L(r)-r$ strongly negative (to $-3.7$\,km), lying outside the
CSR envelope over $r\approx4$--18\,km before returning to $\approx0$ at large $r$. Unlike the elongate Gulf
band, the compact central-basin window produces no spurious large-$r$ positivity. The global envelope test
over $r\le14$\,km gives $p=0.001$ against CSR and $p=0.0025$ against the inhomogeneous null: the regularity
is not an artifact of the basin density gradient.

A counterintuitive sub-result sharpens the interpretation. The complete halokinetic population is regular,
but the mature-diapir sub-class \emph{alone} (\textit{Salzdiapir}, $N=159$) is \emph{clustered} ($R=0.88$;
Table~\ref{tab:R}). Pillows alone are weakly regular ($R=1.11$). Mature diapirs preferentially occupy the
thickest-salt depocentres, so selecting them is a spatially biased sub-sampling that concentrates points and
manufactures clustering; the regularly spaced unit is the full structure population, not a maturity class.

\begin{figure}[t]
\centering
\includegraphics[width=0.92\textwidth]{ngb_type_map.pdf}
\caption{North German salt structures (BGR InSpEE, 693 structures) coloured by internal type. Diapirs,
compressively overprinted diapirs, pillows and double-salinar structures constitute the halokinetic
population; double-salinar structures and many diapirs concentrate in the deep central basin / Gl\"uckstadt
Graben, where regularity is strongest.}
\label{fig:degrmap}
\end{figure}

\begin{figure}[t]
\centering
\includegraphics[width=\textwidth]{ngb_pair_correlation_ripley.pdf}
\caption{Central-basin halokinetic structures ($N=302$). Left: $g(r)$ with CSR and inhomogeneous-Poisson
$95\%$ envelopes (hard-core to $\sim$7\,km, packing peak near 15\,km; global test $p=0.001$ vs.\ CSR,
$p=0.0025$ vs.\ inhomogeneous null). Right: $L(r)-r$ strongly negative (regular) over $r\approx4$--18\,km,
returning to $\approx0$ at large $r$.}
\label{fig:degrgr}
\end{figure}

\subsection{Character of the order: intermediate-range suppression, gradient-limited at long range}
The direct structure factor of the central-basin population (Fig.~\ref{fig:degrsk}) has three features. The
single lowest accessible mode spikes to $S\!\approx\!3.5$; this is the basin-scale density gradient (salt
concentrating in depocentres) and lies within the inhomogeneous-Poisson envelope. Above it, $S(k)$ dips to a
minimum $S_{\min}\!\approx\!0.42$ near $k=0.17\,\mathrm{km}^{-1}$ (a length scale of $\sim$37\,km, roughly
three times the packing spacing), \emph{below} the gradient-matched null --- a genuine suppression of
intermediate-wavelength density fluctuations. The packing peak itself is weak ($S\!\approx\!1.11$ at
$k=0.55\,\mathrm{km}^{-1}$). The number variance is suppressed by $\sim$40\% relative to the
inhomogeneous-Poisson null, but its absolute scaling is super-Poissonian ($\sigma^2\!\sim\!R^{2.8}$,
gradient-dominated, far from the hyperuniform $R^{1}$) and the suppression weakens at the largest $R$.

The character of the salt-structure pattern is therefore a locally rigid, intermediate-range-suppressed
disordered liquid: significantly more ordered than random at short and intermediate range, but \emph{not} a
confirmed hyperuniform pattern. The obstruction to a long-range verdict is physical, not statistical: the
density gradient intrinsic to any salt basin (structures concentrate where salt is thickest) dominates the
lowest-$k$ modes, so $S(k\!\to\!0)\to0$ can be neither confirmed nor excluded from map data alone.

\begin{figure}[t]
\centering
\includegraphics[width=\textwidth]{ngb_structure_factor.pdf}
\caption{Left: direct structure factor $S(k)$ of the central-basin population with CSR and
inhomogeneous-Poisson $95\%$ envelopes. A gradient spike at the lowest mode, a suppression dip
($S\!\approx\!0.42$ at $\sim$37\,km scale) below the gradient-matched null, and a weak packing peak. Right:
number variance $\sigma^2(R)$; suppressed relative to the inhomogeneous-Poisson null but with
gradient-dominated $\sim\!R^{2.8}$ scaling, not the hyperuniform $R^{1}$.}
\label{fig:degrsk}
\end{figure}

\begin{table}[t]
\centering
\caption{Clark--Evans index $R$ (edge-corrected against a matched CSR Monte-Carlo null, $S=499$) and
one-sided dispersion $p$. $R>1$: regular/dispersed; $R\approx1$: random; $R<1$: clustered. NND $=$ median
nearest-neighbour distance.}
\label{tab:R}
\small
\begin{tabular}{lrrrl}
\toprule
Province / subset & $N$ & NND (km) & $R$ & verdict \\
\midrule
\multicolumn{5}{l}{\emph{U.S.\ Gulf Coast (USGS gcdiapirg)}}\\
\quad Whole basin (all pooled)            & 1253 & 11.1 & 0.65 & clustered \\
\quad N.\ Gulf shelf/slope $+$ coast      & 571  & 11.6 & 1.10 & regular ($p=0.002$) \\
\quad Texas--Louisiana shelf core         & 271  & 12.0 & 1.12 & regular ($p=0.002$) \\
\quad Bay of Campeche                     & 366  & 9.2  & 1.02 & random \\
\quad Onshore interior basins (pooled)    & 146  & 13.5 & 0.82 & clustered \\
\midrule
\multicolumn{5}{l}{\emph{North German Basin (BGR InSpEE)}}\\
\quad All structures                      & 693  & 10.1 & 1.22 & regular ($p=0.002$) \\
\quad All halokinetic                     & 571  & 11.1 & 1.21 & regular ($p=0.002$) \\
\quad Central basin, halokinetic          & 302  & 10.3 & 1.41 & regular ($p=0.002$) \\
\quad Gl\"uckstadt Graben, halokinetic    & 109  & 11.0 & 1.47 & regular ($p=0.002$) \\
\quad Salt diapirs only (mature)          & 159  & 11.1 & 0.88 & clustered \\
\quad Salt pillows only                   & 205  & 15.0 & 1.11 & regular ($p=0.004$) \\
\bottomrule
\end{tabular}
\end{table}

\section{Discussion}

\paragraph{A map-view Rayleigh--Taylor signature.}
Within coherent loading-driven provinces, salt-structure populations in both basins are significantly
regularly spaced, with hard-core exclusion and a 12--15\,km characteristic spacing. This is the spatial
expression expected of RT wavelength selection combined with salt-withdrawal exclusion. The hard-core
component has a concrete halokinetic origin: as a structure grows it withdraws salt from its surroundings,
thinning the mother-salt layer into a peripheral rim syncline that may weld out entirely as a primary salt
weld \citep{hudec2007,jackson2017}. Where the source layer is locally exhausted a neighbouring structure
cannot nucleate, so the withdrawal aureole acts as an exclusion radius --- the map-view counterpart of the
observed $g(r)=0$ at small separations --- while the finite RT wavelength sets the dominant spacing between
structures. Two observations argue for an RT origin rather than purely structural control. First, the result
replicates across two independent basins and tectonic settings. Second, regularity strengthens
systematically with salt thickness --- weakest on the Gulf shelf, stronger in the central North German
basin, strongest in the deep, thick-salt Gl\"uckstadt Graben --- consistent with the RT prediction that the
selected wavelength (and hence the clarity of a single dominant spacing) increases with salt-layer
thickness. Purely fault-localised ascent, by contrast, would align and cluster structures, as is indeed seen
where structures are pooled across sub-basins or where ascent follows inherited faults.

\paragraph{Trigger versus spacing control.}
This is not a claim that salt ascent is independent of structural control. Extension and basement faulting
commonly trigger or localise diapirism --- reactive diapirism, in which a fault provides the initial
weakness for salt to exploit \citep{vendeville1992} --- and where ascent remains dominated by inherited
faults the structures align and cluster rather than disperse. Our own data show exactly this: the pooled
onshore interior basins are clustered ($R=0.82$), and salt structures in the fault-influenced parts of the
basins do not contribute the regular signal. The claim here is therefore narrower and explicitly
setting-dependent: \emph{within loading-dominated provinces}, whatever the initial trigger, the
characteristic spacing of the mature structure population is governed by the buoyant (RT) wavelength and
salt-withdrawal exclusion, yielding a statistically regular map pattern; faulting may set where structures
begin, but buoyant wavelength selection sets how far apart they end up. The regularity is thus an emergent
property of the loading-driven provinces, not a universal property of salt tectonics.

\paragraph{Three methodological cautions.}
The analysis exposes three ways to reach a wrong verdict from the same data. (i) \emph{Superposition}:
pooling distinct sub-basins superposes independent processes and manufactures apparent clustering (the
whole-Gulf $R=0.65$, the pooled onshore interior $R=0.82$); inference must be made within single provinces.
(ii) \emph{Sub-class selection}: analysing only mature diapirs returns clustering ($R=0.88$) because
maturity is spatially biased toward depocentres; the regularly spaced unit is the complete structure
population, not a maturity class. (iii) \emph{Density gradient}: the basin-scale concentration of salt in
depocentres injects low-wavenumber power that masks the long-wavelength behaviour and biases the statistics
toward clustering. We control for the first by province subsetting, the second by analysing the full
halokinetic population, and the third by testing against an inhomogeneous-Poisson null --- against which the
detected regularity, being conservative, survives.

\paragraph{Character and limits.}
The two families of statistic deliberately probe different scales: the Clark--Evans index and the
small-$r$ pair correlation measure short-range, hard-core exclusion, whereas the low-wavenumber limit of
$S(k)$ and the large-$R$ number variance measure the long-range suppression of density fluctuations. The
salt-structure populations are unambiguous on the former --- they are significantly short-range regular ---
and gradient-limited on the latter. The structure factor shows that the order extends beyond nearest
neighbours: density fluctuations are
suppressed at intermediate wavelengths ($\sim$3 times the packing spacing) below the gradient-matched null.
But the long-wavelength limit, and hence any claim of hyperuniformity, is inaccessible from map data because
the intrinsic basin density gradient dominates the lowest-$k$ modes. The salt-structure populations are thus
best described as locally rigid, intermediate-range-suppressed disordered liquids. The same point-pattern
and structure-factor diagnostics used here are those applied to disordered, fluctuation-suppressed systems
generally \citep{torquato2018}; their application to salt basins is, to our knowledge, new, and provides a
reproducible, statistically hardened way to quantify halokinetic self-organisation in map view.

\section{Conclusions}
\begin{itemize}\setlength\itemsep{1pt}
\item Salt-structure populations in coherent loading-driven provinces of the U.S.\ Gulf Coast and the North
German Basin are significantly regularly spaced (Clark--Evans $R=1.10$--$1.47$, $p\approx0.002$), with
hard-core exclusion and a 12--15\,km characteristic spacing, robust to a density-gradient
inhomogeneous-Poisson null.
\item Regularity strengthens with salt thickness (Gl\"uckstadt Graben $>$ central basin $>$ Gulf shelf),
consistent with Rayleigh--Taylor wavelength selection plus salt-withdrawal exclusion.
\item The order is real at short and intermediate range, but long-range hyperuniformity cannot be confirmed:
the intrinsic basin density gradient masks the long-wavelength limit.
\item Whole-basin pooling and mature-diapir-only sub-sampling both spuriously return clustering; correct
inference requires single-province analysis of the complete structure population.
\end{itemize}

\section*{Data availability}
U.S.\ Gulf Coast salt diapirs: USGS data release \emph{Salt Diapirs in the Gulf Coast} (gcdiapirg),
\href{https://doi.org/10.5066/P9KH3WB4}{doi:10.5066/\allowbreak P9KH3WB4} (public domain). North German salt structures:
BGR \emph{InSpEE-Salzstrukturen} (\copyright\ BGR, Hannover, 2015), open geodata licence (geonutzv); base
mapping after \citet{reinhold2008}. Analysis used only open scientific Python (NumPy, SciPy, Matplotlib).
All analysis code, the derived coordinate datasets, the figure-generation scripts and the numeric results
are openly available at \url{https://github.com/Ruqing1963/salt-structure-spacing}.

\begin{thebibliography}{99}
\small
\bibitem[Baddeley et al.(2015)]{baddeley2015} Baddeley, A., Rubak, E., Turner, R., 2015.
\textit{Spatial Point Patterns: Methodology and Applications with R.} CRC Press.
\bibitem[BGR(2015)]{bgr} Bundesanstalt f\"ur Geowissenschaften und Rohstoffe, 2015.
InSpEE-Salzstrukturen, \copyright\ BGR, Hannover.
\bibitem[Clark \& Evans(1954)]{clark1954} Clark, P.J., Evans, F.C., 1954. Distance to nearest neighbor as a
measure of spatial relationships in populations. \textit{Ecology} 35, 445--453.
\bibitem[Fernandez \& Kaus(2015)]{fernandez2015} Fernandez, N., Kaus, B.J.P., 2015. Pattern formation in 3-D
numerical models of down-built diapirs initiated by a Rayleigh--Taylor instability.
\textit{Geophys.\ J.\ Int.} 202, 1253--1270.
\bibitem[Hudec \& Jackson(2007)]{hudec2007} Hudec, M.R., Jackson, M.P.A., 2007. Terra infirma: Understanding
salt tectonics. \textit{Earth-Sci.\ Rev.} 82, 1--28.
\bibitem[Jackson \& Hudec(2017)]{jackson2017} Jackson, M.P.A., Hudec, M.R., 2017. \textit{Salt Tectonics:
Principles and Practice.} Cambridge University Press.
\bibitem[Martin(1980)]{martin1980} Martin, R.G., 1980. Distribution of salt structures, Gulf of Mexico.
\textit{U.S.\ Geol.\ Surv.\ Misc.\ Field Studies Map} MF-1213.
\bibitem[Myllym\"aki et al.(2017)]{myllymaki2017} Myllym\"aki, M., Mrkvi\v{c}ka, T., Grabarnik, P., Seijo, H.,
Hahn, U., 2017. Global envelope tests for spatial processes. \textit{J.\ R.\ Stat.\ Soc.\ B} 79, 381--404.
\bibitem[Reinhold et al.(2008)]{reinhold2008} Reinhold, K., Krull, P., Kockel, F., 2008. Salzstrukturen
Norddeutschlands 1:500\,000. Bundesanstalt f\"ur Geowissenschaften und Rohstoffe, Hannover.
\bibitem[Ripley(1977)]{ripley1977} Ripley, B.D., 1977. Modelling spatial patterns. \textit{J.\ R.\ Stat.\
Soc.\ B} 39, 172--212.
\bibitem[Torquato(2018)]{torquato2018} Torquato, S., 2018. Hyperuniform states of matter.
\textit{Phys.\ Rep.} 745, 1--95.
\bibitem[Trusheim(1960)]{trusheim1960} Trusheim, F., 1960. Mechanism of salt migration in northern Germany.
\textit{AAPG Bull.} 44, 1519--1540.
\bibitem[U.S.\ Geological Survey(2021)]{usgsgc} U.S.\ Geological Survey, 2021. Salt Diapirs in the Gulf Coast [gcdiapirg]. U.S.\
Geological Survey data release, doi:10.5066/\allowbreak P9KH3WB4.
\bibitem[Vendeville \& Jackson(1992)]{vendeville1992} Vendeville, B.C., Jackson, M.P.A., 1992. The rise of
diapirs during thin-skinned extension. \textit{Mar.\ Pet.\ Geol.} 9, 331--353.
\end{thebibliography}

\end{document}
